\chapter{评测体系：模型、系统与应用层评测}

\begin{introduction}
\item {\heiti 本章核心问题}：当模型、提示、RAG 和部署都在不断变化时，团队怎样建立一套分层评测体系，知道到底是哪一层在变好、哪一层在拖后腿？
\item {\heiti 主案例推进到哪一步}：为企业知识助手建立第一套可持续运行的评测闭环，让后续优化不再只靠主观感觉。
\item {\heiti 读完后能做什么}：能把评测拆成模型层、系统层和应用层，知道测试集该怎么切、指标该怎么选、RAG 该怎么拆评，以及什么样的分数才足以支撑上线或回滚决策。
\end{introduction}

\section{学习目标}
\begin{itemize}
  \item 理解为什么评测必须分层，而不是只看一个总分或几个示例回答。
  \item 掌握评测集切片、RAG 独立评测、端到端评测、人工评审和线上反馈之间的分工。
  \item 能为企业知识助手设计一套包含正确性、证据性、稳定性和效率的评测矩阵。
  \item 能识别测试集污染、评测泄漏、Judge 偏差和指标错配等常见风险。
\end{itemize}

\section{问题背景}
做到这一步，我们已经有了模型使用、提示工程、微调、RAG 和系统优化的基本工具。问题随之发生了变化：现在团队不再只是“怎么做”，而是“怎么判断做得更好了”。如果没有评测体系，所有优化都只能靠主观印象推进。某次换了提示词后大家觉得回答“似乎更顺了”，某次加了 HyDE 后少数样例“看起来更准了”，但没有人能说清楚究竟改善了哪一类问题、代价是什么、会不会牺牲别的维度。

对企业知识助手来说，这种状态是危险的。因为系统一旦接入真实业务，错误不只是“答得不好看”，还可能意味着制度答偏、权限说错、引用失真、拒答失灵、延迟失控。评测体系的任务，就是把这些模糊的感受翻译成可度量、可复盘、可比较的对象。

上一章我们刚把 RAG 做得更复杂：文档解析、切分、重排、GraphRAG、自适应链路都进来了。系统的活动部件越多，评测就越不能只看一个总分。也正因为如此，本章的重点不是列举一串 benchmark 名称，而是建立一套工程化视角：先决定要评哪一层，再决定用什么数据和什么指标去评，最后决定评测结果如何进入发布和回滚流程。

\section{核心概念}
\subsection{为什么要做分层评测}
分层评测的核心原因很简单：一个企业知识助手的最终回答，是由多层机制共同决定的。模型本身的语言和推理能力是一层，检索、拼接、提示和引用规则是另一层，部署稳定性和延迟又是一层，最后用户在真实业务里是否觉得“能用”，还属于更高一层。

如果只看最终回答分数，团队会失去定位能力。比如答案错误，到底是模型不懂任务、检索没找到证据、提示没有约束住模型、还是系统在高并发下截断了上下文？这些问题在最终结果上都可能表现成“答错了”，但修法完全不同。

因此，评测的第一原则不是“找一个万能总分”，而是先把系统拆层。只有层次拆清楚，优化才有方向，回归才有归因。

\begin{table}[H]
\centering
\small
\begin{tabularx}{\textwidth}{>{\raggedright\arraybackslash}p{2.7cm}>{\raggedright\arraybackslash}p{3.2cm}>{\raggedright\arraybackslash}X}
\toprule
\textbf{评测层} & \textbf{主要问题} & \textbf{典型评测对象} \\
\midrule
模型层 & 模型会不会做这类任务 & 基础问答、抽取、分类、推理、格式遵循 \\
系统层 & 系统能不能把正确材料送进模型并稳定返回 & 检索质量、拼接质量、引用质量、拒答稳定性、延迟 \\
应用层 & 用户在真实业务里是否觉得可用可信 & 任务成功率、人工复核成本、满意度、工单回落率 \\
\bottomrule
\end{tabularx}
\caption{企业知识助手中的三层评测对象}
\label{tab:eval-three-layers}
\end{table}

\subsection{先列出要评哪些能力，再谈怎么打分}
旧稿里有一大块内容在讲“评哪些能力”，新版目前更强调“怎么评”。对教材来说，这两件事都要保留。因为如果团队连自己到底在评什么都没列清楚，就很容易出现“检索评得很细，长上下文和多轮一致性却根本没进表”的情况。

对企业知识助手来说，至少应把几类高频能力显式列出来：理解与抽取、生成与格式遵循、长上下文整合、条件推理、多轮对话一致性，以及诚实地识别知识边界。它们未必都放在同一套 benchmark 里，但应该同时出现在评测资产清单中。

\begin{table}[H]
\centering
\small
\begin{tabularx}{\textwidth}{>{\raggedright\arraybackslash}p{2.9cm}>{\raggedright\arraybackslash}p{3cm}>{\raggedright\arraybackslash}X}
\toprule
\textbf{能力维度} & \textbf{为什么要单列} & \textbf{企业知识助手中的例子} \\
\midrule
理解与抽取 & 决定能否从材料里抓住关键字段和限定条件 & 抽取审批人、报销上限、适用范围 \\
生成与格式遵循 & 决定回答能否按规则交付 & 先给结论、再给引用、按模板输出 \\
长上下文整合 & 决定能否跨段落、跨附件组织答案 & 联合正文、附件和表格解释制度 \\
逻辑推理与条件约束 & 决定是否会漏掉例外条款和前置条件 & “超标可报销，但必须先审批” \\
多轮对话一致性 & 决定澄清、追问和连续问答是否稳定 & 用户追问“那高铁改签呢”时保持上下文一致 \\
诚实与知识边界 & 决定系统会不会在依据不足时乱答 & 文档无答案时拒答，冲突材料时提示人工确认 \\
\bottomrule
\end{tabularx}
\caption{评测不仅要分层，也要明确覆盖哪些能力维度}
\label{tab:eval-capability-matrix}
\end{table}

\subsection{评测集不是一个集合，而是一组切片}
很多团队准备测试集时，会把所有问题丢进一个大表，最后算一个平均分。这种做法在教材和论文里都很常见，但对工程团队来说远远不够。因为平均分会掩盖结构性问题。

例如，同样是 80 分，一个系统可能在 FAQ 上 95 分、在制度条款上 70 分；另一个系统可能在简单问答上 80 分、在高风险权限问答上只有 55 分。这两个 80 分的系统，在业务意义上完全不是一回事。

因此，测试集不应是一个大筐，而应是一组切片。至少要按\textbf{任务类型}、\textbf{风险等级}、\textbf{知识来源}、\textbf{问题表达方式}和\textbf{是否需要引用/拒答}做拆分。切片的好处在于，团队能知道“哪一块真的变好了”，而不是只知道“总体平均值略有波动”。

\begin{table}[H]
\centering
\small
\begin{tabularx}{\textwidth}{>{\raggedright\arraybackslash}p{2.8cm}>{\raggedright\arraybackslash}p{3cm}>{\raggedright\arraybackslash}X}
\toprule
\textbf{切片维度} & \textbf{为什么切} & \textbf{企业知识助手中的例子} \\
\midrule
任务类型 & 不同任务的能力要求不同 & 条款查证、FAQ 问答、流程说明、字段抽取 \\
风险等级 & 高风险问题需要更严格门槛 & 权限、报销、合同、客户承诺问题 \\
知识来源 & 不同来源的质量和结构差别大 & PDF 制度、网页知识库、表格附件、FAQ \\
表达方式 & 口语与规范表述命中率常不同 & “还能报吗” 对比 “超标审批规则” \\
回答约束 & 有的题要求引用，有的题要求拒答 & 有依据问答、证据不足拒答、冲突材料提示 \\
\bottomrule
\end{tabularx}
\caption{测试集应优先具备的切片维度}
\label{tab:eval-slices}
\end{table}

\subsection{公开基准看通用上限，业务金标看真实可用性}
旧稿里提到过 MMLU、C-Eval、CMMLU、AGIEval 等公开基准。教材版里最该保留的，不是把这些名字全背下来，而是先理解它们的角色：\textbf{公开基准更适合看模型底座的通用能力上限，业务金标更适合看你的系统能不能上线。}

公开基准的价值主要有三点。第一，它们可以帮助团队判断某次换底座模型、量化方案或微调路线后，\textbf{通用能力是否被意外打坏}。第二，它们便于和社区常见结果建立大致对照，不至于只在自己的小测试集里自说自话。第三，它们对基础推理、知识面和中文/英文综合能力有一定覆盖，适合做“能力保底线”。

但它们也有非常清楚的边界。企业知识助手最在意的问题，往往是“制度例外有没有漏掉”“权限边界有没有越权”“证据不足时会不会乱答”“回答格式能不能直接进工单系统”。这些问题，公开基准通常并不直接覆盖。因此，一个团队完全可能在 MMLU 或 C-Eval 上分数不错，但在自己最关键的业务场景里依然答得很危险。

\begin{table}[H]
\centering
\small
\begin{tabularx}{\textwidth}{>{\raggedright\arraybackslash}p{2.6cm}>{\raggedright\arraybackslash}p{3.2cm}>{\raggedright\arraybackslash}X}
\toprule
\textbf{评测资产} & \textbf{更适合回答什么} & \textbf{不能替代什么} \\
\midrule
MMLU / C-Eval 一类公开基准 & 模型通用知识、综合理解、基础推理能力是否明显退化 & 企业制度边界、引用规范、拒答策略、权限风险 \\
领域选择题或结构化自动评测集 & 关键知识点、字段抽取、格式遵循是否回归 & 开放式复杂问答中的真实可用性 \\
业务金标开放题集 & 在真实业务问题上是否答对、答稳、答得有依据 & 通用底座的横向可比性 \\
线上回放与灰度样本 & 新版本在真实流量里会不会出事 & 离线阶段的全面错误归因 \\
\bottomrule
\end{tabularx}
\caption{公开基准和业务金标不是二选一，而是职责不同的两层资产}
\label{tab:public-benchmark-vs-business-gold}
\end{table}

因此更稳的评测组合通常是：\textbf{公开基准负责看“底座还在不在”，业务金标负责看“系统还能不能发”。} 两者缺一不可，但绝不能互相替代。

\subsection{自动评测集、人工评测集与保留集要分工}
旧稿里多次强调“自动评测集”和“人工评测集”要并存，这一点在教材里也应该明确写成资产分工。因为评测不只是“样例越多越好”，而是不同资产要承担不同职责。

\begin{table}[H]
\centering
\small
\begin{tabularx}{\textwidth}{>{\raggedright\arraybackslash}p{2.8cm}>{\raggedright\arraybackslash}p{3cm}>{\raggedright\arraybackslash}X}
\toprule
\textbf{资产类型} & \textbf{主要职责} & \textbf{最适合放什么题} \\
\midrule
自动评测集 & 高频回归、快速比较版本、做门禁初筛 & 选择题、字段抽取、格式检查、可脚本核验的问答 \\
人工评测集 & 看业务边界、开放式质量、风险判断 & 高风险条款解释、拒答质量、复杂引用、边界样例 \\
保留集 / 黄金集 & 避免过拟合调参、做最终放行判断 & 团队最在意、但不应天天拿来调试的关键样例 \\
线上坏例回放集 & 把真实失败变成未来回归资产 & 灰度中暴露出的新问法、新资料形态、新错误模式 \\
\bottomrule
\end{tabularx}
\caption{评测资产要像工程资产一样分层治理，而不是所有样例放在一个表里}
\label{tab:eval-asset-roles}
\end{table}

对企业知识助手来说，一个很实用的做法是：自动评测集负责日常回归，人工评测集负责高风险抽检，保留集负责最终放行，线上坏例回放集负责持续生长。这样当团队连续一周都在调 RAG、Prompt 或模型版本时，就不会把所有样例一起“看熟”，也不会在最关键的题上失去最后一道独立验证。

\subsection{评测集也要版本化，不然分数就没有时间坐标}
很多团队会认真记录模型版本、prompt 版本和索引版本，却忽略了另一个同样关键的对象：\textbf{评测集本身也在变。} 线上坏例持续流入、评分 rubric 持续修订、切片边界持续细化，这都意味着“上周的 82 分”和“这周的 82 分”未必是在同一把尺子上量出来的。

因此，评测集最好也像代码一样被版本化管理。这里的“版本”不只是题目列表，还包括切片定义、评分协议、Judge 提示词、人工 rubric 和保留集清单。只要其中任意一项变化了，团队就应该知道自己比较的是“模型变了”，还是“尺子变了”。

\begin{table}[H]
\centering
\small
\begin{tabularx}{\textwidth}{>{\raggedright\arraybackslash}p{3cm}>{\raggedright\arraybackslash}p{3.2cm}>{\raggedright\arraybackslash}X}
\toprule
\textbf{需要版本化的对象} & \textbf{为什么要单独记录} & \textbf{如果不记录会怎样} \\
\midrule
题目与切片清单 & 决定评测覆盖面和难度分布 & 分数变化无法判断是版本进步还是样本换简单了 \\
评分 rubric / Judge 提示 & 决定“什么叫更好”的口径 & 可能把裁判换了误当成模型进步 \\
保留集清单 & 决定最终放行标准是否被偷偷改动 & 最关键的题会在不知不觉中被看熟或被替换 \\
错误标签体系 & 决定坏例统计是否可横向比较 & 同类错误在不同周会被起成不同名字 \\
\bottomrule
\end{tabularx}
\caption{评测集版本化的目的，是让分数有明确时间坐标和口径坐标}
\label{tab:eval-versioning}
\end{table}

教材里把这件事讲清楚很重要，因为工程团队真正要比较的，常常不是“绝对分数有多高”，而是“在同一把尺子下，这次是否真的更稳”。只要评测集不版本化，团队就会在不知不觉中失去这个判断能力。

\subsection{RAG 评测要拆成独立评测与端到端评测}
RAG 系统最容易掉进的评测误区，是只看最终答案。这样做虽然接近用户体验，但无法定位问题来源。更稳妥的做法，是把 RAG 至少拆成两层评：
\begin{itemize}
  \item \textbf{独立评测}：单独评检索层和阅读层，比如相关片段是否被召回、是否进入最终提示、引用是否真实对应来源。
  \item \textbf{端到端评测}：直接评最终答案是否正确、是否有依据、是否在证据不足时正确拒答。
\end{itemize}

独立评测回答的是“证据链是否健康”，端到端评测回答的是“最终行为是否可用”。二者缺一不可。只做独立评测，无法知道模型是否真正用好了证据；只做端到端评测，无法知道系统为什么失败。

\begin{table}[H]
\centering
\small
\begin{tabularx}{\textwidth}{>{\raggedright\arraybackslash}p{2.9cm}>{\raggedright\arraybackslash}p{3.1cm}>{\raggedright\arraybackslash}X}
\toprule
\textbf{评测方式} & \textbf{核心问题} & \textbf{适合回答什么} \\
\midrule
检索独立评测 & 证据有没有被找回来 & 切分、索引、查询增强、重排序是否有效 \\
拼接独立评测 & 正确证据有没有真正入模 & 上下文预算、去重、模板规则是否合理 \\
端到端问答评测 & 用户最终看到的答案是否可靠 & 正确性、引用性、拒答性、整体可用性 \\
\bottomrule
\end{tabularx}
\caption{RAG 评测中独立评测与端到端评测的分工}
\label{tab:eval-rag-levels}
\end{table}

\subsection{指标要围绕“要守住什么”来选}
指标不是越多越好，而是越贴近系统责任越好。对企业知识助手来说，一套可落地的指标通常至少要覆盖四类：
\begin{itemize}
  \item \textbf{正确性}：答案是否事实正确、是否覆盖关键条件。
  \item \textbf{证据性}：答案是否真正由证据支持、引用是否能回到原文。
  \item \textbf{稳定性}：同类问题反复提问时表现是否波动过大。
  \item \textbf{效率性}：延迟、吞吐、超时和失败率是否在可接受范围内。
\end{itemize}

如果系统面向高风险场景，还要额外补一个指标维度：\textbf{拒答质量}。当证据不足时，系统是正确拒答、过度拒答，还是错误自信作答，这三者在业务上完全不同。

\begin{table}[H]
\centering
\small
\begin{tabularx}{\textwidth}{>{\raggedright\arraybackslash}p{2.8cm}>{\raggedright\arraybackslash}p{3.2cm}>{\raggedright\arraybackslash}X}
\toprule
\textbf{指标维度} & \textbf{可选指标} & \textbf{解释} \\
\midrule
检索质量 & Recall@k、MRR、命中率 & 评正确证据能否进入候选前列 \\
答案质量 & 正确率、覆盖率、一致性 & 评最终答案是否答对且答全 \\
证据质量 & 引用准确率、Groundedness & 评答案是否真正受材料支持 \\
拒答质量 & 拒答正确率、误拒率、误答率 & 评系统在证据不足场景下是否稳健 \\
系统质量 & P95 延迟、超时率、失败率 & 评是否具备上线稳定性 \\
\bottomrule
\end{tabularx}
\caption{企业知识助手中常见指标及其职责}
\label{tab:eval-metrics}
\end{table}

\subsection{阈值不是统一门槛：要按风险等级设通过线}
很多团队第一次做评测门禁时，容易直接问一句：“正确率到多少才算能发？”这个问法本身就过于粗糙。因为阈值从来不是一个脱离业务语境的漂亮数字，而是团队对\textbf{失败代价}的显式承诺。产品 FAQ 答错一次，可能只是用户多问一句；报销、权限、合同或合规问题答错一次，代价却可能是审批放错、财务风险或审计事故。两类问题不该共用同一条线。

因此，更稳妥的做法不是先拍一个总阈值，再让所有题往上套，而是先把样例按风险等级分层，再为不同层决定不同的放行逻辑。低风险题可以更看平均表现，中风险题开始关注切片波动，高风险题则常常要看最小通过线、误答率和人工抽检，而不是只看均值。

\begin{table}[H]
\centering
\small
\begin{tabularx}{\textwidth}{>{\raggedright\arraybackslash}p{2.4cm}>{\raggedright\arraybackslash}p{3.1cm}>{\raggedright\arraybackslash}p{3cm}>{\raggedright\arraybackslash}X}
\toprule
\textbf{风险层级} & \textbf{更该优先看的指标} & \textbf{更常见的放行逻辑} & \textbf{企业知识助手中的例子} \\
\midrule
低风险 & 平均正确率、格式稳定性、延迟 & 看总体趋势，允许局部小波动 & 产品 FAQ、通用说明、低风险流程解释 \\
中风险 & 切片正确率、引用准确率、误拒率 & 既看总分，也看关键切片是否退化 & 跨制度问答、跨部门流程说明、表格字段抽取 \\
高风险 & 误答率、拒答质量、人工抽检一致性 & 重点看底线指标和单题抽查，不接受“均值掩盖” & 报销规则、权限边界、合同例外、审批建议 \\
红线场景 & 越权率、伪引用率、审计可追溯性 & 只要触发红线错误就直接阻断发布 & 直接替用户批准付款、捏造制度依据、绕过人工审批 \\
\bottomrule
\end{tabularx}
\caption{不同风险等级下的评测门槛不应共用同一条线}
\label{tab:eval-risk-thresholds}
\end{table}

教材读者真正要学会的，不是记住某个阈值数字，而是记住阈值制定顺序：\textbf{先判失败代价，再挑指标，最后才给门槛。} 如果顺序反过来，团队常会陷入“总分很好看，但高风险样例依旧不敢放”的尴尬状态。

\subsection{诚实性不是附赠项：知识边界识别要单独评}
企业知识助手最危险的失败方式之一，不是“答得短”，而是\textbf{在证据不足时仍然自信作答}。这类风险和普通正确率并不完全重合，所以不能指望它被平均分自动覆盖。旧稿里提到的 Honest 原则，放到这里最值得保留的不是名词，而是一个工程要求：系统必须学会识别自己知道什么、不知道什么，以及当前结论到底有没有足够依据。

因此，诚实性评测最好单独设计样例，而不是只在普通问答里顺便观察。至少要覆盖四类场景：知识库里根本没有答案；材料里只有部分条件；多份文档互相冲突；当前用户无权查看完整证据。对这四类样例，团队真正要看的不是“答得像不像人”，而是系统能否稳定做出保守、透明、可升级的处理。

\subsection{先定位，再打分}
对工程团队来说，一个高分但不可定位的评测体系，价值往往不如一个分数普通但能指出问题层级的体系。因此，评测流程通常应当先做定位，再做打分。一个实用的顺序是：
\begin{enumerate}
  \item 先看错误样例属于哪一层：模型层、检索层、拼接层还是系统层；
  \item 再看它落在哪个切片：口语问法、表格问答、制度例外、高风险拒答等；
  \item 最后才统计分数变化，判断这类问题是否在变多或变少。
\end{enumerate}

这个顺序看起来比“一键跑总分”更麻烦，但它会直接改变团队的优化质量。因为真正推动系统进步的，从来不是一个孤立分数，而是“某类高风险错误正在下降”这种可行动结论。

\subsection{错误分类法要固定，否则坏例很难沉淀成资产}
如果团队每次复盘失败样例时都临时起一个新名字，那么评测集即使越来越厚，也很难变成真正可复用的工程资产。今天说“检索没找准”，明天说“证据不完整”，后天又说“上下文拼坏了”，最后看板上就只会剩一堆无法统计、无法对比的口语化描述。

更成熟的做法是先固定一套足够粗但稳定的错误主类，再允许团队在主类下面补充细粒度备注。主类的目标不是穷尽所有细节，而是让不同版本、不同人、不同周的复盘结果能横向比较。对企业知识助手来说，一套实用的主类通常可以先从下面几类开始：

\begin{table}[H]
\centering
\small
\begin{tabularx}{\textwidth}{>{\raggedright\arraybackslash}p{2.6cm}>{\raggedright\arraybackslash}p{3.3cm}>{\raggedright\arraybackslash}X}
\toprule
\textbf{错误主类} & \textbf{典型表现} & \textbf{第一反应通常该看哪里} \\
\midrule
检索缺失 & 正确证据根本没有进入候选 & 切分、索引、查询改写、召回参数 \\
证据错配 & 引用了材料，但材料并不支持结论 & 引用协议、拼接模板、证据核查规则 \\
推理或抽取漏项 & 证据都在，却漏了条件、例外或字段 & 提示结构、阅读链路、模型能力 \\
边界与拒答错误 & 该拒不拒、该保守不保守、越权给建议 & 风险规则、拒答样本、对齐策略 \\
格式与接口错误 & 结构不符合工单或下游系统要求 & 输出模板、结构化约束、后处理 \\
系统稳定性错误 & 超时、截断、重试失败、并发抖动 & 服务链路、超时策略、资源配置 \\
\bottomrule
\end{tabularx}
\caption{评测错误标签最好先固定主类，再逐步细化子类}
\label{tab:eval-error-taxonomy}
\end{table}

一个简单但很有用的治理动作是：\textbf{每条失败样例至少给一个主标签，必要时再给一个次标签和一个风险标签。} 这样团队既能统计“本周检索缺失有没有下降”，也能继续保留“它发生在高风险权限场景”这种细节。没有标签治理，坏例只会越来越多；有了标签治理，坏例才会慢慢变成可回归、可归因、可训练的资产。

\subsection{人工评测依然重要，但职责要收紧}
大模型时代，自动评测和 LLM-as-a-Judge 很方便，但它们不能完全替代人工评测。尤其在企业问答里，很多高价值判断是业务性的：制度解释是否过度延伸、拒答是否符合风险偏好、引用是否真的足以支持结论、术语是否符合内部习惯。这些内容，自动打分很难完全覆盖。

不过，人工评测也不该被设计成“什么都靠人看”。更好的做法是收紧职责，让人工主要负责三类工作：
\begin{itemize}
  \item 判高风险样例和边界样例；
  \item 校准自动评测与 Judge 结果是否偏移；
  \item 定义新的错误类型并反馈给测试集设计。
\end{itemize}

换句话说，人工评测不是和自动评测竞争，而是给整个评测体系做校准和补盲。

\begin{table}[H]
\centering
\small
\begin{tabularx}{\textwidth}{>{\raggedright\arraybackslash}p{2.9cm}>{\raggedright\arraybackslash}p{3cm}>{\raggedright\arraybackslash}X}
\toprule
\textbf{评测方式} & \textbf{优势} & \textbf{局限} \\
\midrule
规则或脚本评测 & 快、稳、便于回归 & 只能评格式或明确可验证项 \\
LLM-as-a-Judge & 覆盖面广、迭代快 & 会有偏好漂移和提示敏感性 \\
人工评测 & 能看懂业务边界与风险 & 成本高、速度慢、难以全量覆盖 \\
\bottomrule
\end{tabularx}
\caption{三类常见评测方式的能力边界}
\label{tab:eval-methods}
\end{table}

\subsection{人工评测的关键不是多看，而是同口径地看}
人工评测最常见的误区，不是“看得不够多”，而是“每个人看的标准不一样”。如果甲认为“先给结论再给依据”更好，乙认为“必须先逐条解释再给结论”更好，丙又把“更礼貌”看得比“更保守”重要，那么你最后得到的并不是团队偏好，而只是几位评审各自的写作习惯。

因此，人工评测真正的第一任务不是扩样本，而是把口径磨平。一个可执行的最小流程通常包括四步：先拿一小批代表性样例做双人标注；再集中复盘分歧最大的题；把分歧背后的优先级写回 rubric；最后用修订后的 rubric 重跑一轮校准。只有这四步跑通了，扩大人工评测规模才值得。

对工程团队来说，最该警惕的信号不是“分数高低”，而是“分歧集中在同一类题”。如果权限场景、拒答场景或跨文档证据场景总是出现相反判断，说明问题不在评审员不认真，而在协议没有写清楚。此时继续堆更多人工评测，只会更快地产生更多互相冲突的标签。

\subsection{LLM-as-a-Judge 要想用得稳，先把评分协议写清楚}
旧稿里讲过 PandaLM、对话竞技场一类自动化评测思路。放到教材里，最重要的不是记住具体平台名，而是理解：\textbf{Judge 评测不是“让另一个模型随便打分”，而是要先定义评分协议。}

对工程团队来说，至少有三种常见 Judge 用法：
\begin{itemize}
  \item \textbf{规则化单点评分}：按“是否给结论、是否有引用、是否越权延伸”这类 rubric 打分；
  \item \textbf{成对比较}：让 Judge 在 A/B 两个回答之间选更好者，适合版本对比；
  \item \textbf{竞技场式偏好汇总}：收集多轮 pairwise 结果，形成更接近 Chatbot Arena 的偏好排序。
\end{itemize}

\begin{table}[H]
\centering
\small
\begin{tabularx}{\textwidth}{>{\raggedright\arraybackslash}p{3cm}>{\raggedright\arraybackslash}p{3cm}>{\raggedright\arraybackslash}X}
\toprule
\textbf{Judge 方式} & \textbf{更适合什么任务} & \textbf{主要提醒} \\
\midrule
Rubric 单点评分 & 结构化问答、拒答规范、引用完整性 & rubric 不清时，分数会飘得很厉害 \\
Pairwise 对比 & 新旧版本回答谁更好 & 容易受输出顺序和措辞偏好影响，需要换序复测 \\
Arena 式汇总 & 多版本长期比较与排序 & 成本更高，且需要更强的样本覆盖与统计稳定性 \\
\bottomrule
\end{tabularx}
\caption{Judge 评测的关键不是“用没用模型”，而是协议是否稳定}
\label{tab:judge-modes}
\end{table}

无论用哪一种，Judge 都最好配三道护栏：第一，保留一小批人工金标样本，用来校准 Judge 是否漂移；第二，对高风险题做顺序打乱或双向比较，避免“先看到的回答”天然占优；第三，Judge 提示词与评分 rubric 要版本化，否则你看到的“模型进步”可能只是裁判口径变了。

\subsection{工具、赛制与平台只是载体，不是评测本身}
旧稿里提过 Chatbot Arena、OpenAI Evals、PandaLM 这类工具或赛制。放到新版教材里，更应该把它们理解成三种载体，而不是三种“自动得到真相”的答案。Evals 一类框架适合做模板化回归和版本记录；Judge 模型适合放大人工评审能力；Arena 或 pairwise 赛制适合比较多个版本之间的偏好差异。

但无论工具怎么选，它们都替代不了三件基础工作：先切清任务和风险样例；再准备足够稳定的金标或参考协议；最后把结果接进发布门禁。如果这三件事没做好，换任何平台都只是在把混乱自动化。

\subsection{RAGAS、ARES 一类自动评测框架能加速，但别让框架替你定义目标}
旧稿里专门讲过 RAGAS 和 ARES，这些内容不该在新版里只剩延伸阅读。对工程团队来说，它们真正的价值是：\textbf{帮你更快地把 RAG 评测跑起来}，尤其是在检索相关性、答案忠实度、上下文利用等维度上做大规模回归。

但它们的边界同样要讲清楚。第一，框架里的指标定义通常仍然依赖提示词、Judge 或少量人工标注校准，因此它们不是绝对真理。第二，它们更擅长做\textbf{批量趋势判断}，不擅长替你拍板高风险业务结论。第三，如果团队连“什么叫好的企业回答”都没定义清楚，那么把框架接进来只是更快地产生一堆不知该不该信的分数。

因此，RAGAS、ARES 一类框架最适合放在下面这个位置：先有小规模人工金标和清楚的评分协议，再用自动评测框架把同类口径放大到更大样本；最后仍然保留高风险题的人类抽检与发布门禁。这样它们会成为效率工具，而不是判断外包工具。

\subsection{评测成本也要分层：把最贵的人力留给最危险的样例}
评测不仅要追求准确，还要考虑成本。很多团队一开始就想“所有题都人工看一遍”，很快就会被速度和人力拖垮；另一种极端是“所有题都交给自动评测”，结果又在高风险场景里失去判断力。更现实的做法，是像管理算力预算一样管理评测预算。

\begin{table}[H]
\centering
\small
\begin{tabularx}{\textwidth}{>{\raggedright\arraybackslash}p{2.7cm}>{\raggedright\arraybackslash}p{2.8cm}>{\raggedright\arraybackslash}p{2.8cm}>{\raggedright\arraybackslash}X}
\toprule
\textbf{评测层} & \textbf{单位成本} & \textbf{最适合承担什么职责} & \textbf{什么时候不能只靠它} \\
\midrule
脚本与规则评测 & 最低 & 高频回归、格式检查、接口约束、红线扫描 & 需要判断业务边界或证据充分性时 \\
Judge 批量评测 & 中等 & 趋势比较、版本 A/B、RAG 批量回归 & 高风险结论需要最终拍板时 \\
人工评测 & 最高 & 高风险抽检、协议校准、分歧复盘、门禁兜底 & 样本量很大、只需要日常回归时 \\
\bottomrule
\end{tabularx}
\caption{评测预算也应分层使用，而不是让最贵的手段承担全部工作}
\label{tab:eval-cost-layers}
\end{table}

这张表背后的教材化结论是：\textbf{把最贵的人力，留给最有否决权的样例。} 低风险、高频、结构化任务优先自动化；趋势比较交给 Judge 放大；真正决定“发不发”的高风险切片，再由人工做最终把关。这样评测体系才既可持续，又不会在关键处失守。

\subsection{测试集污染和评测泄漏会让分数失真}
当团队越来越频繁地根据固定样例调提示、调检索和调模型时，一个很现实的风险会越来越大：测试集开始被“看熟了”。这就是测试集污染。更严重时，测试集内容甚至直接流入训练、微调、提示模板或 RAG 素材，导致评测分数失真。

对企业知识助手来说，最常见的泄漏路径包括：把评测题直接拿去做 few-shot 示例；把失败样例原封不动加入知识库；在调试时不断对着同一批测试题微调检索器；把评测集答案写进运维 FAQ 却不更新测试口径。

因此，评测集必须像数据资产一样被治理：区分开发集、回归集、保留集；记录每次评测所用的版本；高价值保留集不要天天拿来调参；一旦测试集进入训练或规则设计，就需要重新补充新题。

\subsection{评测的终点是发布门禁，而不是报告}
很多团队能写出很漂亮的评测报告，却没有把报告变成发布动作的一部分。这样一来，评测就只是“看一眼”，而不是“卡一下”。真正成熟的做法，是让评测结果进入上线门禁：
\begin{itemize}
  \item 高风险切片低于阈值，不能发；
  \item 拒答误答率上升，不能发；
  \item P95 延迟超预算，不能发；
  \item 总分略涨但关键切片明显下滑，也不能发。
\end{itemize}

只要评测不进入门禁，团队就很容易被少数漂亮样例诱导，做出“整体看起来不错”的错误发布决策。

\subsection{门禁也要分层：阻断项、预警项和观察项别混在一起}
很多团队明明已经有评测，却依然做不出稳定的发布决策，一个常见原因就是：\textbf{所有指标都叫“重要”，却没有说清楚谁真的能拦版本。} 结果往往是看板上同时摆着二三十个数字，大家开会时却只能凭感觉争论“这次到底算不算能发”。更稳的做法，是把门禁指标至少分成三类。

\begin{table}[H]
\centering
\small
\begin{tabularx}{\textwidth}{>{\raggedright\arraybackslash}p{2.6cm}>{\raggedright\arraybackslash}p{3cm}>{\raggedright\arraybackslash}X}
\toprule
\textbf{门禁类型} & \textbf{主要职责} & \textbf{企业知识助手里常放什么} \\
\midrule
阻断项 & 只要触发就不能发 & 高风险误答率、越权率、伪引用率、系统失败率红线 \\
预警项 & 允许进入灰度，但必须被明确关注 & 某些中风险切片波动、Judge 与人工分歧上升、P95 接近预算 \\
观察项 & 不直接阻断，但帮助长期判断趋势 & 风格偏好、回复长度变化、低风险 FAQ 小幅波动 \\
\bottomrule
\end{tabularx}
\caption{发布门禁不是“所有指标一票否决”，而是要先区分阻断、预警和观察}
\label{tab:gate-tiers}
\end{table}

这样分层之后，团队的沟通会清楚很多。阻断项回答“能不能发”，预警项回答“发了以后最该盯什么”，观察项回答“系统长期往哪个方向漂”。如果这三类不拆开，团队就很容易要么被一堆小波动吓得不敢发，要么被平均分掩盖掉真正该拦的风险。

\subsection{发布门禁必须绑定回滚动作和责任人}
仅仅知道“哪些指标能拦版本”还不够。真正进入发布现场后，团队还必须继续回答两个问题：\textbf{谁有权决定回滚，以及触发后具体怎么回滚。} 如果门禁只是停留在看板展示层，那么一旦线上灰度出现高风险误答，大家仍然可能围着图表讨论十分钟，最后错过最该果断回退的窗口。

\begin{table}[H]
\centering
\small
\begin{tabularx}{\textwidth}{>{\raggedright\arraybackslash}p{2.9cm}>{\raggedright\arraybackslash}p{3cm}>{\raggedright\arraybackslash}p{2.8cm}>{\raggedright\arraybackslash}X}
\toprule
\textbf{触发信号} & \textbf{默认动作} & \textbf{谁至少要被拉进来} & \textbf{为什么要提前写清} \\
\midrule
阻断项离线不过线 & 版本不得进入灰度 & 模型负责人、发布负责人 & 避免“先发一点看看”式的侥幸发布 \\
灰度中高风险误答率突增 & 立即暂停放量或回退上一稳定版 & 值班负责人、业务 owner、平台负责人 & 高风险错误扩散速度常比会商速度更快 \\
Judge / 人工分歧突然增大 & 进入人工复核，不继续自动放量 & 评测负责人、业务评审 & 这通常说明口径变了，不能按旧门禁继续推 \\
延迟或失败率超预算 & 回退流量、降级链路或切回旧模型 & 平台负责人、服务值班 & 系统层退化会抵消所有离线收益 \\
\bottomrule
\end{tabularx}
\caption{门禁不只是“能不能发”，还应提前定义“出事后谁来按哪一个按钮”}
\label{tab:gate-actions}
\end{table}

这张表背后的教材化结论是：\textbf{发布机制必须把判断、动作和责任绑在一起。} 只有这样，评测才不是一份被阅读的报告，而是一套会真正改变发布行为的控制系统。

\subsection{离线高分不等于线上安全：灰度、A/B 与样本回放}
旧稿还专门强调过综合评测平台和发布后的持续验证。这一点在教材里也应该明确保留：\textbf{离线评测解决的是“值不值得发”，线上灰度解决的是“发出去后会不会出事”。}

对企业知识助手来说，更稳的发布顺序通常是：
\begin{enumerate}
  \item 先跑离线回归，确认关键切片不过线就不发布；
  \item 再做小流量灰度或影子流量回放，观察真实查询分布下的行为；
  \item 必要时做 A/B，对比新旧版本在拒答率、人工复核率、P95 延迟和用户满意度上的变化；
  \item 最后才逐步放量，而不是一次性全量替换。
\end{enumerate}

\begin{table}[H]
\centering
\small
\begin{tabularx}{\textwidth}{>{\raggedright\arraybackslash}p{3cm}>{\raggedright\arraybackslash}p{3.2cm}>{\raggedright\arraybackslash}X}
\toprule
\textbf{线上验证动作} & \textbf{主要看什么} & \textbf{为什么离不开它} \\
\midrule
影子流量回放 & 同一请求在新旧版本上的输出差异 & 能提前发现离线没覆盖到的真实流量问题 \\
小流量灰度 & 错答、拒答、超时、人工接管率 & 能控制风险，把坏影响限制在小范围内 \\
A/B 对比 & 用户反馈、完成率、复核率、延迟 & 能避免只看主观个例做版本判断 \\
发布后回放集更新 & 新增线上坏例子是否进入回归集 & 能让评测体系跟着真实问题一起成长 \\
\bottomrule
\end{tabularx}
\caption{评测闭环不在报告结束，而在发布后继续生长}
\label{tab:online-eval-actions}
\end{table}

\subsection{Judge 自己也要校准：裁判漂了，分数就不再可比}
LLM-as-a-Judge 很好用，但它还有一个常被忽略的现实：\textbf{Judge 本身也是模型，也会漂。} 当 Judge 模型版本更新、提示协议改变、评分 rubric 微调，甚至只是上下文样式变了，分数都可能整体偏移。若团队没有单独校准，就会把“裁判口径变了”误判成“被评系统进步了”。

因此，Judge 最好也像业务模型一样有自己的小回归集。这个回归集不需要很大，但应覆盖三类样例：一批团队已经高度共识的高质量回答，一批明显错误或越界的低质量回答，以及一批最容易引发分歧的边界题。每次 Judge 模型、Judge prompt 或评分协议变动时，都先在这批样例上跑一轮，确认排序和分数没有整体漂到不可接受的程度，再把它投入大规模评测。

这条纪律尤其重要，因为企业知识助手后面常常会把 Judge 接进发布流程。只要 Judge 结果开始影响“能不能发”，它就不再只是一个方便的小工具，而是一个带着决策影响力的评审组件。带影响力的组件，就必须自己接受校准。

\section{主案例推进}
在企业知识助手的主案例里，评测体系的第一版不是追求庞大，而是追求\textbf{够用且能驱动发布决策}。我们先把问题切成四组：制度条款查证、权限流程说明、产品 FAQ、证据不足拒答。每组再按高风险和普通风险分层。

\subsection{先做一套能复用的样例结构}
为了让后续每一章的优化都能复用，我们给每条评测样例统一保存以下字段：
\begin{itemize}
  \item 原始问题；
  \item 所属切片；
  \item 标准答案或参考判断；
  \item 允许引用的标准来源；
  \item 是否应该拒答；
  \item 失败时的错误类型标签。
\end{itemize}

有了这套结构，团队以后做 RAG 优化、提示调整、模型切换和部署压测时，就能持续复用同一套评测资产，而不是每次从头再整理一批零散样例。

\subsection{把评测结果做成“分层看板”}
在这个阶段，我们不追求一个复杂平台，先做一张简单但够用的评测看板：
\begin{enumerate}
  \item 模型层看：基础问答和格式遵循是否回归；
  \item 系统层看：检索命中率、引用准确率、拒答误答率、P95 延迟；
  \item 应用层看：高频任务成功率、人工复核率和用户反馈。
\end{enumerate}

这样做的价值在于，当团队引入 HyDE、重排序或新的模型版本时，能直接看到改变主要落在哪一层，而不是只看一个混合后的总分。

\begin{table}[H]
\centering
\small
\begin{tabularx}{\textwidth}{>{\raggedright\arraybackslash}p{2.8cm}>{\raggedright\arraybackslash}p{3cm}>{\raggedright\arraybackslash}p{3cm}>{\raggedright\arraybackslash}X}
\toprule
\textbf{看板层} & \textbf{核心指标} & \textbf{主要用途} & \textbf{谁最关心} \\
\midrule
模型层 & 任务正确率、格式遵循率 & 判断模型或提示是否回归 & 算法与应用开发 \\
系统层 & Recall@k、引用准确率、误答率、延迟 & 判断 RAG 和服务链路是否健康 & 平台与后端工程 \\
应用层 & 高风险任务成功率、人工复核率、满意度 & 判断业务是否真正受益 & 产品、业务与运营 \\
\bottomrule
\end{tabularx}
\caption{企业知识助手中的分层评测看板}
\label{tab:eval-dashboard}
\end{table}

\subsection{把评测结果变成发布规则}
在主案例中，团队最终约定了三条最硬的上线规则：
\begin{enumerate}
  \item 高风险切片不能只看平均分，必须逐项抽查并满足最小阈值；
  \item 证据不足场景的误答率不能上升；
  \item 系统层延迟和失败率不能为了换来一点准确率而明显恶化。
\end{enumerate}

这三条规则看起来保守，但它们保证了企业知识助手不会因为一次“答得更像人”的改动，就在高风险场景里答得更危险。

\subsection{把灰度坏例直接流回回放集，不要只留在群聊里}
主案例推进到这一步时，团队通常会遇到一个特别真实的现象：真正有价值的新坏例，往往不是离线设计出来的，而是灰度期间由真实用户问出来的。问题在于，这些坏例如果只停留在群聊截图、临时工单或口头复盘里，几周后就会再次出现，团队却只能重复惊讶。

因此，企业知识助手的评测闭环里最好明确规定：\textbf{所有进入灰度复盘会的高价值坏例，都要在会后 24 小时内变成回放样例。} 至少补齐原始问题、真实输出、期望行为、所属切片、错误标签和风险等级，然后进入回放集或保留集。这样一来，线上问题才会真正反哺离线评测，而不是每次都以“新问题”的名义重新出现。

从教材视角看，这一步特别值得强调，因为它把评测从“静态题库”推进成了“会生长的资产系统”。评测体系真正成熟的标志，不是第一次做得多完整，而是它能不能持续吸收真实世界暴露出来的新边界。

\subsection{总分上涨但关键切片下滑时，默认按失败处理}
主案例里最难的一类发布判断，不是“整体都变差了”，而是“平均分变好了，但高风险切片变差了”。这类情况最容易诱导团队自我说服：总分既然上升了，是不是可以先发一小版看看？教材里必须把结论说得更硬一点：\textbf{只要关键切片承担的是业务否决权，总分上涨也不能覆盖它的下滑。}

\begin{table}[H]
\centering
\small
\begin{tabularx}{\textwidth}{>{\raggedright\arraybackslash}p{3.2cm}>{\raggedright\arraybackslash}p{3cm}>{\raggedright\arraybackslash}X}
\toprule
\textbf{评测结果组合} & \textbf{默认动作} & \textbf{原因} \\
\midrule
总分上升，关键高风险切片也上升 & 可以进入灰度候选 & 收益与风险方向一致 \\
总分上升，但关键高风险切片下滑 & 默认不发布 & 平均收益不能覆盖关键风险恶化 \\
总分持平，但误答率明显下降 & 可作为保守升级候选 & 在高风险场景里，降误答常比涨均分更有价值 \\
总分上升，但延迟和失败率超预算 & 默认不发布 & 系统层退化会把离线收益抵消为线上事故 \\
\bottomrule
\end{tabularx}
\caption{发布判断要让关键切片拥有真正的否决权}
\label{tab:eval-release-decision}
\end{table}

这也是为什么成熟团队会在门禁里明写“哪几项拥有 veto 权”。一旦 veto 条件清楚，评测就不再只是复盘材料，而会真正变成版本决策机制。

\section{常见坑与工程取舍}
\begin{itemize}
  \item \textbf{只看总分，不看切片}：平均分会掩盖高风险问题的恶化。
  \item \textbf{只做端到端评测}：出了问题时无法知道是检索错、拼接错还是模型错。
  \item \textbf{测试集长期不更新}：团队越调越像在刷榜，分数越来越高，但真实效果未必更好。
  \item \textbf{把 Judge 分数当成真理}：Judge 方便，但它也会偏，尤其对业务边界和拒答质量不一定敏感。
  \item \textbf{评测不进入发布门禁}：没有门禁的评测，最后往往只剩报告，没有约束力。
\end{itemize}

还需要补四个很现实的取舍：
\begin{itemize}
  \item \textbf{先做大而全测试集还是先做关键切片}：工程起步期通常先做关键切片。
  \item \textbf{先追求自动化还是先追求标签质量}：高风险场景里，标签质量通常更重要。
  \item \textbf{先补更多指标还是先补更强归因}：能定位问题层级的评测，往往比多几个漂亮指标更有用。
  \item \textbf{先拉高正确率还是先压低误答率}：在企业知识助手里，高风险场景通常先压误答率。
\end{itemize}

\section{本章练习}
\begin{exercise}[指标设计]
如果企业知识助手要回答报销制度、权限审批和产品 FAQ 三类问题，你会至少保留哪三类指标来支撑发布决策？请说明每类指标在守住什么风险。
\end{exercise}

\begin{solution}
至少应保留三类指标。第一类是{\heiti 正确性指标}，例如答案正确率或关键条件覆盖率，用来守住“有没有答偏”。第二类是{\heiti 证据性指标}，例如引用准确率或 groundedness，用来守住“答案是否真由材料支持”。第三类是{\heiti 风险与稳定性指标}，例如误答率、误拒率、P95 延迟或失败率，用来守住“高风险场景是否在乱答，以及系统是否能稳定提供服务”。如果是高风险企业场景，这三类指标通常缺一不可。
\end{solution}

\begin{exercise}[测试集污染]
为什么团队不能长期对着同一批测试题不断调提示、调检索和调模型？如果必须频繁回归测试，应该怎样降低测试集污染风险？
\end{exercise}

\begin{solution}
因为长期对着同一批题反复调参，会让系统越来越“熟悉”这批题，分数上升却不代表泛化能力真的变强，这就是测试集污染。降低风险的做法包括：区分开发集、回归集和保留集；高价值保留集不要天天拿来调参；记录每次评测所用版本；当测试样例已经实质进入提示、训练或规则设计时，及时补充新题并更新保留集。
\end{solution}

\section{延伸阅读}
\begin{itemize}
  \item \textbf{RAGAS、ARES 等 RAG 评测框架资料}：适合理解检索与答案联合评测的常见实现思路。
  \item \textbf{各类模型评测与 LLM-as-a-Judge 讨论资料}：适合理解自动评测为什么高效，也为什么需要人工校准。
  \item \textbf{回归测试与发布门禁工程实践}：适合理解评测结果如何真正进入上线流程，而不是停留在报告层。
  \item \textbf{数据集治理与测试集污染相关资料}：适合理解为什么测试资产也需要版本化和生命周期管理。
\end{itemize}

\section{小结}
评测体系的核心，不是把系统打出一个漂亮总分，而是让团队知道：哪一层在变好、哪一层在退化、哪些错误可以接受、哪些错误必须拦住。对企业知识助手来说，真正有用的评测一定是分层的、切片的、可回归的，并且最终会进入发布门禁，而不是只停留在一份好看的评测报告里。

\section{下一章过渡}
当团队能够稳定看见“系统在变好还是变坏”之后，下一个问题就会变成：如果我们想继续把模型行为推向更安全、更符合人类偏好，应该怎么做？下一章将进入对齐与强化学习，讨论为什么需要 RLHF、PPO 在其中扮演什么角色，以及这条路线到底适合解决什么问题。
